\documentclass[letterpaper,11pt]{article}

\usepackage{latexsym}
\usepackage[empty]{fullpage}
\usepackage{titlesec}
\usepackage{marvosym}
\usepackage[usenames,dvipsnames]{color}
\usepackage{verbatim}
\usepackage{enumitem}
\usepackage[hidelinks]{hyperref}
\usepackage{fancyhdr}
\usepackage[english,russian]{babel}
\usepackage{tabularx}
\input{glyphtounicode}

%----------FONT OPTIONS----------
% Поддержка кириллицы и современный шрифт без засечек
\usepackage[T2A]{fontenc}
\usepackage[utf8]{inputenc}
\usepackage[sfdefault]{roboto} % Если шрифт не находится, замените на \usepackage[default]{sourcesanspro}
\renewcommand{\familydefault}{\sfdefault}

\pagestyle{fancy}
\fancyhf{} % clear all header and footer fields
\fancyfoot{}
\renewcommand{\headrulewidth}{0pt}
\renewcommand{\footrulewidth}{0pt}

% Adjust margins
\addtolength{\oddsidemargin}{-0.5in}
\addtolength{\evensidemargin}{-0.5in}
\addtolength{\textwidth}{1in}
\addtolength{\topmargin}{-.5in}
\addtolength{\textheight}{1.0in}

\urlstyle{same}

\raggedbottom
\raggedright
\setlength{\tabcolsep}{0in}

% Sections formatting
\titleformat{\section}{
  \vspace{-4pt}\scshape\raggedright\large
}{}{0em}{}[\color{black}\titlerule \vspace{-5pt}]

% Ensure that generate pdf is machine readable/ATS parsable
\pdfgentounicode=1

%-------------------------
% Custom commands
\newcommand{\resumeItem}[1]{
  \item\small{
    {#1 \vspace{-2pt}}
  }
}

\newcommand{\resumeSubheading}[4]{
  \vspace{-2pt}\item
    \begin{tabular*}{0.97\textwidth}[t]{l@{\extracolsep{\fill}}r}
      \textbf{#1} & #2 \\
      \textit{\small#3} & \textit{\small #4} \\
    \end{tabular*}\vspace{-7pt}
}

\newcommand{\resumeSubSubheading}[2]{
    \item
    \begin{tabular*}{0.97\textwidth}{l@{\extracolsep{\fill}}r}
      \textit{\small#1} & \textit{\small #2} \\
    \end{tabular*}\vspace{-7pt}
}

\newcommand{\resumeProjectHeading}[2]{
    \item
    \begin{tabular*}{0.97\textwidth}{l@{\extracolsep{\fill}}r}
      \small#1 & #2 \\
    \end{tabular*}\vspace{-7pt}
}

\newcommand{\resumeSubItem}[1]{\resumeItem{#1}\vspace{-4pt}}

\renewcommand\labelitemii{$\vcenter{\hbox{\tiny$\bullet$}}$}

\newcommand{\resumeSubHeadingListStart}{\begin{itemize}[leftmargin=0.15in, label={}]}
\newcommand{\resumeSubHeadingListEnd}{\end{itemize}}
\newcommand{\resumeItemListStart}{\begin{itemize}}
\newcommand{\resumeItemListEnd}{\end{itemize}\vspace{-5pt}}

%-------------------------------------------
%%%%%%  RESUME STARTS HERE  %%%%%%%%%%%%%%%%%%%%%%%%%%%%

\begin{document}

%----------HEADING----------
\begin{center}
    \textbf{\Huge Аслан Зейналов} \\ \vspace{4pt}
    \small Аналитик данных (Data Analyst) \\ \vspace{4pt}
    \small +7 (927) 230-33-31 $|$ \href{mailto:aslanzeinalov02@gmail.com}{aslanzeinalov02@gmail.com} $|$ 
    \href{https://t.me/aslanzei}{@aslanzei} $|$
    \href{https://github.com/AslanZeinalov}{GitHub} $|$ Казань, 24 года
\end{center}

%-----------ABOUT-----------
\section{О себе}
В процессе работы над учебными и практическими проектами приобрел уверенные навыки работы с данными: от сбора и очистки до поиска инсайтов и построения прогнозных моделей. Умею находить закономерности и аномалии в данных, используя их для улучшения бизнес-показателей. В данный момент нахожусь в поиске компании для прохождения стажировки и трудоустройства. С примерами моих проектов и исследований можно ознакомиться в профиле на GitHub.

%-----------EDUCATION-----------
\section{Образование}
  \resumeSubHeadingListStart
    \resumeSubheading
      {Школа 21 (Сбер)}
      {} % <-- Добавлен город (4-й аргумент)
      {Программа Data Science}
      {Октябрь 2025 -- настоящее время} % 
      \resumeItemListStart
        \resumeItem{Прошел конкурсный отбор на программу обучения.}
      \resumeItemListEnd
      
    \resumeSubheading
      {Уфимский топливно-энергетический колледж}
      {} % <-- Добавлен город (4-й аргумент)
      {Технология бурения нефтяных и газовых скважин}
      {Сентябрь 2018 -- Июнь 2022}
      \resumeItemListStart
        \resumeItem{Окончил с отличием (красный диплом).}
      \resumeItemListEnd
  \resumeSubHeadingListEnd

%-----------EXPERIENCE-----------
\section{Опыт работы}
  \resumeSubHeadingListStart
    \resumeSubheading
      {ООО «ВБ-ВОСТОК»}
      {} % <-- Добавлен город (4-й аргумент)
      {Специалист отдела инвентаризации $\rightarrow$ Помощник руководителя отдела}
      {Август 2022 -- Октябрь 2026} % <-- Исправлен год с 2026 на настоящее время
      \resumeItemListStart
        \resumeItem{Составлял отчеты в Excel по заполнению мест хранения, учету потерь и проведению работ на участке (сводные таблицы, формулы).}
        \resumeItem{Управлял штатом из 5 сотрудников, занимался постановкой и контролем выполнения задач.}
      \resumeItemListEnd
  \resumeSubHeadingListEnd

%-----------PROJECTS-----------
\section{Проекты}
    \resumeSubHeadingListStart
    
      \resumeProjectHeading
          {\textbf{A/B-тест формы оплаты} $|$ \emph{Python, pandas, scipy, EDA, статистические тесты}}
          {} % <-- Добавлен 2-й аргумент (год или ссылка)
          \resumeItemListStart
            \resumeItem{Провёл полный цикл продуктового анализа: от очистки данных и EDA до статистической проверки гипотезы. Выявил прирост конверсии +14.6 п.п. (p-value = $4 \cdot 10^{-14}$) и подготовил рекомендацию по запуску новой формы на 100\% пользователей.}
            \resumeItem{\href{https://github.com/AslanZeinalov/AB-test-analysis}{Ссылка на репозиторий (GitHub)}}
          \resumeItemListEnd
          
    \resumeProjectHeading
          {\textbf{Комплексный анализ продаж и логистики сети веломагазинов} $|$ \emph{Python, SQL, EDA, бизнес-аналитика}}
          {} % <-- Добавлен 2-й аргумент (год или ссылка)
          \resumeItemListStart
            \resumeItem{Обнаружил упущенную выручку в \$1.37 млн из-за логистических проблем и предложил решения по оптимизации запасов.}
            \resumeItem{Провел RFM-сегментацию: выделил 26.4\% VIP-клиентов, приносящих 47.7\% выручки, и разработал стратегию их удержания.}
            \resumeItem{\href{https://github.com/AslanZeinalov/retail_store_analysis}{Ссылка на репозиторий (GitHub)}}
          \resumeItemListEnd
          
      \resumeProjectHeading
          {\textbf{Дашборд для анализа интернет-магазина (138K+ заказов)} $|$ \emph{Power BI, DAX, Power Query, ETL}}
          {} % <-- Добавлен 2-й аргумент
          \resumeItemListStart
            \resumeItem{Создал 6 интерактивных дашбордов с продвинутыми DAX-мерами (YoY growth, LTV, Repeat Rate) для мгновенной оценки бизнес-метрик, для CEO и руководителей отделов}
            \resumeItem{Спроектировал модель данных (Star Schema) и настроил ETL-процессы для автоматического обновления отчетов.}
            \resumeItem{\href{https://github.com/AslanZeinalov/E-Commerce-Sales-Analytics-dashboard}{Ссылка на репозиторий (GitHub)}}
          \resumeItemListEnd
          
    \resumeSubHeadingListEnd

%-----------SKILLS-----------
\section{Навыки}
 \begin{itemize}[leftmargin=0.15in, label={}]
    \small{\item{
     \textbf{Hard Skills:} Python, SQL, PostgreSQL, Power BI, DAX, Power Query, Excel, pandas, numpy, matplotlib, scipy, scikit-learn, Machine Learning, Jupyter, Google Colab, EDA, AB-тесты, статистические тесты \\
     \textbf{Soft Skills:} Презентация проектов, объяснение сложного простым языком, работа в команде, быстрая обучаемость \\
     \textbf{Языки:} Русский (Родной), Английский (B1 -- Intermediate)
    }}
 \end{itemize}

%-------------------------------------------
\end{document}
